\begin{corollary}[Discriminant and level in the rank-three count]\label{rk3:volume}
Let \(D=|\operatorname{disc}K|\) for an imaginary quadratic field \(K\).
Let \(S\) consist of the two classical lists of
Proposition~\ref{rk3:count}, at principal levels \(N\) over \(K\)
and \(ND\) over \(\mathbf Q\), respectively, with one entry per
system. For integers \(N,e\ge1\),
\[
 |S|\ll h_K D^4 N^{18}e^6+(ND)^4e
       \ll D^5N^{18}e^6 .
\]
The constants are absolute. Consequently the integer of
Lemma~\ref{rk3:integer}, under its trace hypotheses, satisfies
\[
 \log|C_p|\ll |J|D^5N^{18}e^7\log p .
\]
\end{corollary}
\begin{proof}
Put \(L=3N\) and \(X=\mathrm{SL}_3(\mathbf C)/\mathrm{SU}(3)\)
with a fixed invariant Riemannian metric. As in
Proposition~\ref{rk3:count}, pass to principal level \(L\) and
divide out the real split centre. At infinity
\(\mathrm{GL}_3(\mathbf C)/(\mathbf R_{>0}\mathrm U(3))\simeq X\).
Determinant and strong approximation for \(\mathrm{SL}_3\) identify
the connected components with the finite ray class group modulo
\(L\mathcal O_K\). In fact the determinant image of the principal
compact open is the subgroup of finite integral idelic units
congruent to one at that modulus, as is seen on diagonal matrices.
The map to the ordinary class group has kernel a quotient of
\((\mathcal O_K/L\mathcal O_K)^\times\). There are thus at
most \(h_K L^2\) components.

For a finite adelic representative \(g\) of a component put
\(\Lambda_v=g_v\mathcal O_{K_v}^3\). A member of its arithmetic
stabilizer preserves these lattices and acts trivially modulo \(L\).
Its determinant is a unit \(\epsilon\in\mathcal O_K^\times\) with
\(\epsilon\equiv1\bmod L\). If \(\epsilon\ne1\), then
\(0<|\operatorname{N}_{K/\mathbf Q}(\epsilon-1)|\le4\), whereas
divisibility by \(3\) makes that norm divisible by \(9\).
Hence the stabilizer is the level-\(L\) subgroup of
\[
 \Gamma_g^1=\mathrm{SL}_3(K)\cap
          \prod_{v<\infty}\operatorname{Stab}_{\mathrm{SL}_3(K_v)}
                                     (\Lambda_v).
\]
Its index in \(\Gamma_g^1\) is at most \(L^{16}\). Indeed, in a
basis of each \(\Lambda_v\) the image of reduction at
\(\mathfrak p_v^b\Vert L\mathcal O_K\) has size at most
\[
 |\mathrm{SL}_3(\mathcal O_{K_v}/\mathfrak p_v^b)|
 =q_v^{8(b-1)}|\mathrm{SL}_3(\mathbf F_{q_v})|\le q_v^{8b}.
\]
Reduction is surjective, since a lift with unit determinant can be
corrected by scaling one row. Each successive kernel is the additive
space of trace-zero matrices over the residue field. Multiplication
over \(v\) gives \(\operatorname{N}(L\mathcal O_K)^8=L^{16}\).

Prasad's covolume formula applies to the split simply connected
\(K\)-group \(\mathrm{SL}_3\). For every \(g\) each local stabilizer
is hyperspecial, with integral model \(\mathrm{SL}(\Lambda_v)\)
and reductive special fibre \(\mathrm{SL}_3\). They form a coherent
collection, equal to the standard one almost everywhere. Strong
approximation at the noncompact complex place gives the principal
arithmetic group appearing in the formula. Its Haar covolume in
\(\mathrm{SL}_3(\mathbf C)\) is
\[
 c_\infty D^{8/2}\prod_{v<\infty}
   \frac{q_v^8}{|\mathrm{SL}_3(\mathbf F_{q_v})|}
 =c_\infty D^4\zeta_K(2)\zeta_K(3).
\]
Here the Tamagawa number is one and the archimedean constant is
fixed for the single complex place and the chosen Haar measure.
In particular there is no dependence on the ideal class of \(g\).
The local zeta factors give \(\zeta_K(j)\le\zeta(j)^2\) for \(j>1\).
Prasad,
\emph{Volumes of \(S\)-arithmetic quotients of semi-simple groups},
\S\S3.3, 3.7 and 3.11~\cite{prasadVolume}.

The level-\(L\) subgroup acts faithfully and without torsion, by
its action on the underlying integral lattice modulo three. Its
quotient of \(X\) is noncompact since \(\mathrm{SL}_3\) is \(K\)-split
of positive rank. Its Riemannian volume is at most a fixed constant
times \(D^4L^{16}\), by the covolume and index bounds. Gelander's
theorem gives a homotopy-equivalent simplicial complex with at most
a fixed multiple of this volume in vertices and uniformly bounded
vertex degree. A vertex of degree \(b\) belongs to at most \(2^b\)
simplices, so the same linear bound holds for all simplices.
Local-coefficient cochains contribute at most \(e^6\) dimensions per
simplex. Summing over the \(h_K L^2\) possible components and
using cuspidal cohomology as in Proposition~\ref{rk3:count} proves
the first term. The second follows from the same proposition over
the fixed field \(\mathbf Q\) in rank two with \(M=ND\) and
coefficient dimension \(2e\).
Gelander,
\emph{Homotopy type and volume of locally symmetric manifolds},
Theorem 1.5(1)~\cite{gelanderVolume}.

Minkowski supplies each ideal class with an integral ideal of norm
at most \(\sqrt D\). There are at most \(\tau(m)\) ideals of norm
\(m\), by the quadratic prime splitting possibilities. Consequently
\(h_K\le\sqrt D(1+\log\sqrt D)\le D\).
Finally \(q\le p^2\), so the per-factor logarithmic cost
\(2e\log q+1\) of Lemma~\ref{rk3:integer} is \(O(e\log p)\).
\end{proof}
